% !TeX root = ../main.tex
\chapter{Algebras, representations and states}\label{f:quantum}
The bracket in Chapter~\ref{f:symplectic} gives a number for each pair of linear observable labels. It does not give the expectation of either observable, much less the expectation of their product. To see what the bracket does determine, start with a real vector space $V$ and an alternating form $\sigma$. For the scalar, $v=[f]$ and $\sigma([f],[h])=\langle f,Eh\rangle$, not $\omega(Ef,Eh)$. We shall use this one input to construct two quantum algebras. A state is an extra choice; a Hilbert space can then be constructed rather than postulated. We display $\hbar>0$ in the defining relations and specify when natural units $\hbar=1$ are used.

\section{An algebra by generators and relations}
A complex unital algebra is a complex vector space with an associative bilinear multiplication and an identity $1$. An involution $*$ is conjugate-linear and satisfies $(AB)^*=B^*A^*$, $(A^*)^*=A$ and $1^*=1$. It models taking an adjoint; it is defined even before there are operators to take adjoints of. The elements $A=A^*$ are called Hermitian. Imposing a relation $R=0$ means also imposing $ARB=0$ for arbitrary $A,B$. Accordingly one divides by a \emph{two-sided ideal}, a linear subspace stable under multiplication on both sides. For an involution on the quotient one also needs $I^*=I$. These conditions make $[A][B]=[AB]$ and $[A]^*=[A^*]$ independent of representatives. Any unital $*$-homomorphism killing $I$ factors uniquely through the quotient, by sending $[A]$ to the image of $A$.

\begin{definition}[Polynomial canonical commutation relations]\label{f:polynomial-ccr}
Complexify $V$ to $V_{\mathbb C}$. Using the \hyperref[f:tensor-product]{algebraic tensor product}, form the tensor algebra $T(V_{\mathbb C})=\bigoplus_{n\ge0}V_{\mathbb C}^{\otimes n}$, with degree zero $\mathbb C$, product concatenation, and star conjugating entries and reversing their order. The polynomial CCR algebra is its quotient by the two-sided $*$-ideal generated, for $v,w\in V$, by
\[
 v\otimes w-w\otimes v-i\hbar\sigma(v,w)1 .
\]
The adjoint of each displayed relation is its negative. The classes $\Phi(v)$ of degree-one tensors are real-linear in $v$, are Hermitian for real $v$, and satisfy $[\Phi(v),\Phi(w)]=i\hbar\sigma(v,w)1$, where $[A,B]=AB-BA$.
\end{definition}
These are finite words and their finite sums, not convergent power series. The algebra is \emph{universal}: given Hermitian elements $T_v$ of another unital $*$-algebra with this linearity and commutator, multiply the $T_v$ in the order of each word and extend linearly. The generating relations map to zero, hence so does their ideal, and the map on the quotient is the unique homomorphism with $\Phi(v)\mapsto T_v$. The relations are consistent: on the vector space of commuting polynomials $\operatorname{Sym}(V_{\mathbb C})$ let $M_v$ multiply by $v$ and let the derivation $D_v$ send $w$ to $\sigma(v,w)$. Then $[D_v,M_w]=\sigma(v,w)1$; hence $M_v+(i\hbar/2)D_v$ obey the desired commutators and the unit acts nontrivially. This last check only uses an associative algebra, not a Hilbert-space inner product.

If $r\in\operatorname{rad}\sigma$, $\Phi(r)$ commutes with every word. Conversely a central linear generator has $\sigma(r,v)=0$ for every $v$. One may retain these central quantities or prescribe their values in a sector, for example $\Phi(r)=\lambda(r)1$ for a specified real-linear $\lambda$. Imposing every classical condition as an ideal of an unreduced quantum algebra would be different: in a single pair with $[Q,P]=i\hbar 1$, the two-sided ideal generated by $Q$ already contains $1$. Its quotient describes no canonical pair at all.

\section{Bounded Weyl generators and the meaning of regularity}
A \emph{representation} assigns to each algebra element $A$ an operator $\pi(A)$, linearly, with $\pi(1)$ the identity and $\pi(AB)=\pi(A)\pi(B)$. A Hilbert-space representation of the polynomial algebra generally needs \emph{unbounded} operators. Specify a common dense domain $\mathcal D_\pi$ carried into itself by all the operators, so every stated product acts there. To represent the involution we require $\langle\xi,\pi(A)\eta\rangle=\langle\pi(A^*)\xi,\eta\rangle$ on this domain. A Hermitian symbol acts symmetrically, not necessarily self-adjointly after closure. The Weyl construction avoids these domain questions in the algebra multiplication itself. A $C^*$-algebra is a complete normed $*$-algebra with submultiplicative norm and $\|A^*A\|=\|A\|^2$; bounded operators on a Hilbert space are the model \cite{BR}.

\begin{definition}[Weyl relations]\label{f:weyl-relations}
Use one unitary symbol $W(v)$ for each $v\in V$, with
\begin{equation}
 W(0)=1,\qquad W(v)^*=W(-v),\qquad
 W(v)W(w)=e^{-i\hbar\sigma(v,w)/2}W(v+w).
 \label{f:weyl-phase}
\end{equation}
The Weyl algebra $\mathcal W(V,\sigma)$ is the universal unital $C^*$-algebra with these relations. The exponent convention is $W(tv)=\exp(it\Phi(v))$ when a field operator exists; there is no further factor of $1/\hbar$ in this convention.
\end{definition}
To construct it, first take finite formal sums of the symbols. The two parenthesizations of three factors have the same phase because $\sigma(u,v)+\sigma(u+v,w)=\sigma(v,w)+\sigma(u,v+w)$. Reversing the product conjugates its phase, as required by $*$. Take the supremum of the operator norm over Hilbert-space realizations by unitaries: the norm of $\sum c_j W(v_j)$ is at most $\sum |c_j|$. It separates elements. Indeed on $\ell^2(V)$ with orthonormal vectors $e_w$ indexed by the labels, put
\[
 U(v)e_w=e^{-i\hbar\sigma(v,w)/2}e_{v+w}.
\]
These operators permute a basis up to phases, obey \eqref{f:weyl-phase}, and a finite sum with distinct labels applied to $e_0$ records every coefficient. Completion in the supremum norm gives the $C^*$-algebra. Each unitary realization extends by continuity, uniquely since the finite span is dense. Also $W(v)W(tw)=e^{-i\hbar t\sigma(v,w)}W(tw)W(v)$. Commutation for all real $t$ and all $w$ is equivalent to $v\in\operatorname{rad}\sigma$: in particular nondegeneracy excludes nonzero central Weyl labels, not the scalar unit.

For example, to remove the radical $R$ in its zero-charge sector, divide by the closed ideal generated by $W(r)-1$, $r\in R$. In this quotient $W(v+r)=W(v)W(r)$ has the same class as $W(v)$. These classes obey the Weyl relations of $(V/R,\bar\sigma)$; conversely sending $W(v)$ to the generator labeled by $v+R$ kills the ideal. The two maps are inverse on generators. Retaining a central label, fixing a character on it, and discarding the label without specifying a sector are different operations.

A representation is \emph{regular} if $t\mapsto\pi(W(tv))$ is strongly continuous for each $v$, that is, continuous on each Hilbert vector in norm. Stone's theorem then gives a self-adjoint generator $\widehat\Phi(v)$ for this one-parameter group. One may differentiate the Weyl relation in both parameters on common smooth vectors for a finite collection of directions, obtaining $[\widehat\Phi(v),\widehat\Phi(w)]=i\hbar\sigma(v,w)1$. Such vectors are dense: for the finite-dimensional span $L$, $(t,v)\mapsto e^{it}\pi(W(v))$ is a continuous unitary representation of $\mathbb R\times L$ with group law $(t,v)(s,w)=(t+s-\hbar\sigma(v,w)/2,v+w)$. Averaging a vector against smooth, compactly supported approximate identities on this group gives smooth approximants. Differentiating is not an operation in the abstract $C^*$-norm. The realization $U$ above is not even regular when $V\ne0$: for $v\ne0$ and $t\ne0$, $\|U(tv)e_0-e_0\|=\sqrt2$. Nor do symmetric polynomial generators satisfying a commutator on a domain, without self-adjointness and control of exponentiation, give the Weyl relations. The two universal constructions answer different questions.

\section{Positivity constructs a cyclic representation}
A \emph{state} on a unital $*$-algebra $\mathcal A$ is a complex-linear functional $\rho$ with $\rho(1)=1$ and $\rho(A^*A)\ge0$ for every $A$. Positivity is stronger than assigning real numbers to Hermitian elements: it constrains measurements of sums of observables. In particular $(A,B)_\rho=\rho(A^*B)$ is a positive semidefinite Hermitian form (conjugate-linear in the first entry). Polarizing its real nonnegative diagonal gives Hermitian symmetry; positivity on $A+zB$ for all complex $z$ gives $|(A,B)_\rho|^2\le(A,A)_\rho(B,B)_\rho$. Consequently a vector of zero length pairs to zero with every vector.

\begin{theorem}[Gelfand--Naimark--Segal construction]\label{f:gns}
Every state is the vector state of a cyclic representation. On the domains generated by the cyclic vectors this realization is unique up to a unitary taking one such vector to the other. For a general $*$-algebra these are common invariant dense domains; for a $C^*$-algebra the representation extends by bounded operators to the completed Hilbert space.
\end{theorem}
\begin{proof}
Set $N_\rho=\{A:(A,A)_\rho=0\}$. By the inequality this is the subspace orthogonal to everything. If $A\in N_\rho$ then $\rho((BA)^*BA)=(A,B^*BA)_\rho=0$ for every $B$, so it is a \emph{left} ideal. On the quotient vector space define
\[
 \langle[A],[B]\rangle=\rho(A^*B),\qquad
 \pi_\rho(B)[A]=[BA],\qquad \Omega_\rho=[1].
\]
The left ideal is exactly what makes left multiplication well defined here; we do not multiply two classes as elements of a quotient algebra. The inner product is positive definite, and the adjoint pairing holds since $(BA)^*C=A^*B^*C$. Complete to $\mathcal H_\rho$. The vectors $\pi_\rho(A)\Omega_\rho=[A]$ form the dense domain (this is cyclicity) and $\langle\Omega_\rho,\pi_\rho(A)\Omega_\rho\rangle=\rho(A)$. In the $C^*$ case, $\|B\|^2 1-B^*B$ has a positive square root, so positivity on its sandwich between $A^*$ and $A$ gives $\|[BA]\|^2\le\|B\|^2\|[A]\|^2$. Thus each operator extends to the completion. Finally, in any cyclic realization of the same state, $[A]\mapsto\pi(A)\Omega$ preserves these inner products and has dense range. Its unitary extension intertwines on the generated domain, and everywhere for bounded operators.
\end{proof}

A real-linear bijection $S:V_1\to V_2$ preserving the forms gives a $*$-isomorphism $\alpha_S$ by $W_1(v)\mapsto W_2(Sv)$; apply universality in both directions. The same argument applies to polynomial generators. A state $\rho_2$ pulls back to $\rho_2\circ\alpha_S$. For two \emph{specified} states, a unitary intertwiner taking one cyclic vector to the other exists precisely when $\rho_1=\rho_2\circ\alpha_S$, by the inner-product argument above. Another question is whether an automorphism is implemented in a \emph{fixed} representation by $\pi(\alpha_S(A))=U\pi(A)U^*$. Such a unitary need not fix the state vector; it need not exist at all (Exercise~\ref{f:ex-implementation}).

\section{Characteristic functions and mixed covariances}
A Weyl state is determined by its characteristic function $v\mapsto\rho(W(v))$, since finite sums are dense. For a real symmetric bilinear form $\mu$, consider the centered Gaussian candidate $\chi(v)=\exp[-\mu(v,v)/2]$. Which $\mu$ define states? The answer includes a restriction on the correlations between different labels, not just on each variance.

\begin{proposition}[Gaussian positivity]\label{f:gaussian-positivity}
The candidate is a state characteristic function if and only if $\Lambda=\mu+i\hbar\sigma/2$, extended sesquilinearly to $V_{\mathbb C}$, is positive semidefinite. The Gaussian representation is regular, with ordered two-point function $\langle\Omega_\rho,\widehat\Phi(v)\widehat\Phi(w)\Omega_\rho\rangle=\Lambda(v,w)$.
\end{proposition}
\begin{proof}
Positivity on finite Weyl sums is exactly positive semidefiniteness of the matrix $K_{jk}=\rho(W(v_j)^*W(v_k))$ for every finite list. For the candidate,
\[
 K_{jk}=e^{-\mu(v_j,v_j)/2}\,e^{\Lambda(v_j,v_k)}\,e^{-\mu(v_k,v_k)/2} .
\]
If $[\Lambda(v_j,v_k)]$ is a Gram matrix of vectors, its entrywise $n$th power is a Gram matrix of their $n$-fold tensor powers. Summing with weights $1/n!$, including the all-ones matrix at $n=0$, proves positivity of the entrywise exponential. The outside diagonal factors preserve it. The resulting state on finite sums extends: its cyclic representation has unitary Weyl generators, hence is bounded by the defining supremum norm. Conversely the Gram matrices of $(W(tv_j)-1)/t$ in the state are positive. Their entries tend to $\Lambda(v_j,v_k)$: expand $K_{jk}$ at labels $tv_j,tv_k$, subtract $\chi(-tv_j)$ and $\chi(tv_k)$, and add $1$ before dividing by $t^2$. Limits preserve positivity.

On any finite-dimensional label span, the kernel $\langle\pi(W(x))\Omega_\rho,\pi(W(y))\Omega_\rho\rangle$ is smooth in $(x,y)$ by the displayed formula. The vector orbit of $\Omega_\rho$ is smooth there as well: the inner products of two difference quotients are mixed difference quotients of this kernel, so those vector quotients are Cauchy in norm. The same argument applies to iterated differences. Weyl translates of $\Omega_\rho$ have the same property and span a dense subspace; boundedness of the unitaries gives strong continuity everywhere. The ordered product in the proposition is thus taken on a smooth vector, not in the abstract $C^*$-algebra. Applying $-\partial_s\partial_t$ to $e^{-i\hbar st\sigma(v,w)/2}\chi(sv+tw)$ at zero gives its expectation.
\end{proof}
In particular $\mu(v,v)\mu(w,w)\ge\mu(v,w)^2+\hbar^2\sigma(v,w)^2/4$. The commutator fixes the imaginary part of the ordered moment but not the real covariance $\mu(v,w)$ \cite{WaldQFT}.

\begin{corollary}[Ordered Wick pairings]\label{f:wick-pairings}
In the centered Gaussian state above, also called a centered \emph{quasifree} state, take real labels $v_1,\ldots,v_n$. Write $X_j=\widehat\Phi(v_j)$ and $\Lambda_{jk}=\Lambda(v_j,v_k)$. The cyclic-vector expectation in the specified operator order is
\[
 \langle X_1\cdots X_n\rangle_\rho=
 \begin{cases}
  0, & n\text{ odd},\\[2pt]
  \displaystyle\sum_{\mathcal P}\prod_{(j,k)\in\mathcal P}\Lambda_{jk}, & n\text{ even}.
 \end{cases}
\]
Here $\mathcal P$ runs over partitions into pairs, writing the earlier index first in each pair. No time ordering is understood.
\end{corollary}
\begin{proof}
Put $u=\sum_j t_jv_j$ with real parameters. Multiplying the Weyl operators in the written order, each new factor acquires a phase from every earlier factor. Hence their generating function is
\begin{align*}
 Z(t)&=\rho\big(W(t_1v_1)\cdots W(t_nv_n)\big)
 =e^{-(i\hbar/2)\sum_{j<k}t_jt_k\sigma(v_j,v_k)}\chi(u)\\
 &=\exp\!\left[-\tfrac12\sum_j t_j^2\mu(v_j,v_j)
                       -\sum_{j<k}t_jt_k\Lambda_{jk}\right].
\end{align*}
On the smooth cyclic vector of Proposition~\ref{f:gaussian-positivity} the ordered moment is $i^{-n}\partial_{t_1}\cdots\partial_{t_n}Z(0)$. The diagonal terms cannot contribute to the coefficient of $t_1\cdots t_n$. Every contributing term selects disjoint pairs covering all the indices, so $n=2m$. The $m!$ possible orders of the selected terms cancel the exponential's $1/m!$; their sign $(-1)^m$ cancels $i^{-2m}$. This proves the formula. Regularity alone gives field generators, not moments of every order in an arbitrary state vector; even when those moments exist, they need not obey this Gaussian rule.
\end{proof}

For $X=\widehat\Phi(v)$ and $Y=\widehat\Phi(w)$ the four-field example is
\[
 \langle X^2Y^2\rangle_\rho
 =\Lambda(v,v)\Lambda(w,w)+2\Lambda(v,w)^2 .
\]
Subtracting the mean from each square leaves two cross pairings, both with $X$ before $Y$: the last factor is a square, not an absolute square. More generally $N(z)=e^{\mu(z,z)/2}\pi(W(z))$ generates \emph{Wick powers} by $i^{-k}\partial_s^k N(sz)|_{s=0}$ on smooth vectors. Removing the contractions internal to each factor leaves the cross term in the identity
\[
 N(sv)N(tw)=e^{-st\Lambda(v,w)}N(sv+tw),
\]
obtained directly from the Weyl relation. In differentiating three times in each variable, two cross contractions choose two out of three slots on each side and a bijection between those slots: $\binom32^2\,2!=18$. The remaining derivatives leave one field from each side in a single Wick product. This is the multiplicity in the fish term \eqref{f:ext-fish}; time ordering and extension to coincident points are separate operations there. Here all scalar field labels are smeared tests modulo the equation of motion.

\begin{example}[An oscillator in position space]\label{f:oscillator-state}
For $v=(a,b)$, $w=(c,d)$ let $\sigma(v,w)=ad-bc$. On $L^2(\mathbb R,\dd x)$ take $Q\psi=x\psi$, $P\psi=-i\hbar\psi'$ on Schwartz functions, so that $\widehat\Phi(v)=aQ+bP$ there. Its Weyl group acts by
\[
 (W(a,b)\psi)(x)=e^{ia(x+\hbar b/2)}\psi(x+\hbar b).
\]
Composing two such translations gives the phase difference $\hbar(bc-ad)/2$ relative to $W(a+c,b+d)$, just as in \eqref{f:weyl-phase}. Take the normalized oscillator wave function $\psi_0(x)=(\pi l^2)^{-1/4}e^{-x^2/(2l^2)}$, $l>0$. In the expectation of $W(a,b)$ the product of Gaussians is centered at $x=-\hbar b/2$. Shifting to $y=x+\hbar b/2$ leaves the Fourier integral of $e^{-y^2/l^2}$, giving
\[
 \langle\psi_0,W(a,b)\psi_0\rangle
   =\exp\!\left[-\tfrac14(l^2 a^2+\hbar^2 b^2/l^2)\right].
\]
The two variances are $l^2/2$ and $\hbar^2/(2l^2)$; their product saturates the bound. Neither the width nor the vector was specified by the Weyl relations.
\end{example}

Two copies already distinguish symplectic and statistical independence. Set $\hbar=l=1$ here. Translate the first oscillator wave function by a random displacement $X$ of mean zero and Gaussian variance $s^2$, and the second by $\varepsilon X$ with $\varepsilon=+1$ or $-1$. Average these product vector states; convex averaging preserves positivity. Translation of a wave function to center $X$ multiplies its characteristic function by $e^{iaX}$. The joint characteristic functions are therefore
\[
 \chi_\varepsilon=\exp\!\left[-\tfrac14\sum_{j=1}^2(a_j^2+b_j^2)
                          -\tfrac{s^2}{2}(a_1+\varepsilon a_2)^2\right].
\]
Setting either pair of labels to zero gives identical marginals for the two signs. Yet $\rho(Q_1Q_2)=\varepsilon s^2$ while $[Q_1,Q_2]=0$. These correlations need not even be entanglement: the mixture was constructed from product vectors. For two general label subspaces, take finite lists from each. If $C$ is the trial real cross-covariance matrix, Gaussian positivity requires
\[
 \begin{pmatrix}
  \mu_1+i\hbar\sigma_{11}/2 & C+i\hbar\sigma_{12}/2\\[2pt]
  (C+i\hbar\sigma_{12}/2)^* & \mu_2+i\hbar\sigma_{22}/2
 \end{pmatrix}\ \ge 0
\]
for \emph{every} choice of lists. Here $*$ is conjugate transpose, $\mu_j$ are the two marginal blocks, and $\sigma_{jk}$ are fixed by the joint bracket. Any linear combination representing a zero label must pair to zero with both lists. A tensor product of marginal states does not in general solve this extension problem. For a spacetime field, a further question is whether the two-point function is a distribution with admissible singularities (Chapter~\ref{f:microlocal}).

\section{Which generators belong to the whole?}
There is a situation where a tensor product is justified at the algebraic level. If $V=V_1\oplus V_2$ and $\sigma(V_1,V_2)=0$, the Weyl generators of the summands commute and $W(v_1+v_2)=W(v_1)W(v_2)$. Conversely any two Weyl realizations with commuting ranges give a realization of this direct sum by that product. This is the universal property of the maximal $C^*$ tensor product (the completion allowing all commuting representations). It neither requires nor selects a product state. If the mixed bracket is nonzero, operators on separate tensor factors would commute and give the wrong phase. If labels are shared, treating the copies as independent misses their identification. Even when two subalgebras commute, this does not establish a factorization $\mathcal H=\mathcal H_1\otimes\mathcal H_2$ in which each local algebra acts as all bounded operators on its own Hilbert factor \cite{Haag}.

There is also a generation question even before a state is specified. For subspaces $V_1,V_2\subset V$ the Weyl relations generate the labels in $V_1+V_2$ regardless of their mixed bracket. No other single Weyl label is in the generated $C^*$-subalgebra: in the basis representation above, its elements send $e_0$ into the closed span of $e_v$ with $v\in V_1+V_2$, whereas $W(w)e_0=e_w$ is orthogonal to that span if $w$ is outside the sum. Checking the bracket of known generators does not test whether a proposed observable is generated.

For a nonlinear example return to the two arcs of Example~\ref{f:two-rotor}. In the algebra of configuration functions take the finite Laurent polynomials in the two open transports $z_j=e^{i\alpha_j}$. At the shared endpoints the relative frame rotation sends $(z_1,z_2)$ to $(tz_1,t^{-1}z_2)$, $|t|=1$. The monomial $z_1^k z_2^l$ has weight $k-l$; linear independence of the Laurent monomials makes the invariant subalgebra exactly the Laurent polynomials in $w=z_1z_2$. This $U(1)$ Wilson loop gives the multiplication operator $(U\psi)(\alpha)=e^{i\alpha}\psi(\alpha)$ on $L^2(S^1)$. With angular momentum $L=-i\hbar\partial_\alpha$, differentiating the product gives $[L,U]=\hbar U$ on smooth periodic functions, in agreement with $\{w,e\}=iw$. If one retains on each arc only functions already invariant under independent changes of its endpoint frames, each arc algebra contains just constants within this configuration class. Their products cannot distinguish $w=1$ from $w=-1$. The lost observable is recovered by multiplying the \emph{open} transports and then taking invariants, not by changing a state or a commutator.

\section{Local operator algebras and a separating collar}\label{f:local-factors}
A bounded-operator algebra in a representation has another natural closure besides norm closure. For a set $S\subset\mathcal B(\mathcal H)$ its \emph{commutant} $S'$ consists of all bounded operators commuting with every element of $S$. The double commutant theorem says that for a unital $*$-algebra $S$, $S''$ is its closure for weak operator convergence (convergence of every matrix coefficient). A unital $*$-algebra with $\mathfrak M=\mathfrak M''$ is a von Neumann algebra; it is a \emph{factor} when its center $\mathfrak M\cap\mathfrak M'$ consists only of scalars \cite{BR}.

The Hilbert tensor product is the completion of the algebraic tensor product for the inner product obtained by extending $\langle x\otimes y,x'\otimes y'\rangle=\langle x,x'\rangle\langle y,y'\rangle$ sesquilinearly. For example, take $\mathfrak M=M_2(\mathbb C)\otimes 1$ on $\mathbb C^2\otimes\mathbb C^n$. In a two-by-two array of $n$-dimensional blocks, an operator commuting with $E_{11}\otimes1$ is block diagonal. Commuting also with $E_{12}\otimes1$ forces the diagonal blocks to agree. Conversely $1\otimes T$ commutes with the whole first algebra, so $\mathfrak M'=1\otimes M_n(\mathbb C)$ and the center is scalar. For $n=2$ the Bell vector $(e_1\otimes e_1+e_2\otimes e_2)/\sqrt2$ gives probability $1/2$ for finding the first basis state on both sides, not the product $1/4$ of the marginal probabilities. The state is entangled; the tensor factors still exist.

Projections distinguish a failure of factorization that is independent of this correlation test. A projection $p=p^*=p^2$ is \emph{minimal} in $\mathfrak M$ if it is nonzero and has no nonzero proper subprojection in $\mathfrak M$. In the matrix example $E_{11}\otimes1$ is minimal in $\mathfrak M$ even though its range has dimension $n$. A factor with a minimal projection is type~I; in a suitable Hilbert tensor decomposition it acts as $\mathcal B(\mathcal K)\otimes1$. A projection is \emph{finite} if no proper subprojection $q<p$ is equivalent to it, where equivalence means $v^*v=p$, $vv^*=q$ for some partial isometry $v\in\mathfrak M$. A factor is type~III precisely when it has no nonzero finite projection; in particular it has no minimal projection. The absence of minimal projections alone does not characterize type~III.

Consider the vacuum representation $\pi_0$ of the massive real free scalar in four-dimensional Minkowski space. For a nonempty bounded double cone $O=I^+(p)\cap I^-(q)$, let $\mathfrak R(O)$ be the double commutant of the Weyl unitaries $\pi_0(W([f]))$ with $f$ supported in $O$. A structural theorem for this field says that $\mathfrak R(O)$ is a type~III factor \cite{Haag,Yngvason}. It cannot be the full bounded-operator algebra on a Hilbert factor associated to the sharp region: such an algebra contains minimal projections. This conclusion uses the local field and its representation, not just a pair of nonzero vacuum correlations. Restricting the vacuum to the region means evaluating $\langle\Omega,A\Omega\rangle$ for $A\in\mathfrak R(O)$; it requires no partial trace.

If the closure of a double cone $O_1$ lies inside a larger double cone $O_2$, the two boundaries leave a collar. The \emph{split property} of this field \cite{Haag,FewSplit} supplies a type~I factor $\mathfrak N$ with
\[
 \mathfrak R(O_1)\subset\mathfrak N\subset\mathfrak R(O_2).
\]
In a tensor decomposition for $\mathfrak N$, the inner algebra acts on the first factor, while $\mathfrak R(O_2)'\subset\mathfrak N'$ acts on the second. The latter commutant contains observables localized spacelike to $O_2$, by locality. It is the \emph{intermediate} algebra that is type~I, not either double-cone algebra. The split theorem does not supply a limiting type~I factor when the collar is removed.

\section*{Exercises}
\begin{exercise}[Basic: null vectors are not relations]
Take $\mathcal A=M_2(\mathbb C)$, $\rho(A)=\langle e_1,Ae_1\rangle$ for the first coordinate vector. Identify $N_\rho$, exhibit its failure to be two-sided, and identify the cyclic Hilbert space.
\end{exercise}
\noindent\emph{Solution.} $\rho(A^*A)=\|Ae_1\|^2$, so $N_\rho$ consists of matrices with zero first column. For matrix units, $E_{12}\in N_\rho$ but $E_{12}E_{21}=E_{11}\notin N_\rho$. The isometry $[A]\mapsto Ae_1$ is onto $\mathbb C^2$, and the representation is ordinary matrix multiplication on vectors.

\begin{exercise}[Intermediate: why no bounded polynomial pair?]
Suppose bounded operators on a nonzero Hilbert space satisfy $[A,B]=1$. Use $[A,B^n]$ to obtain a contradiction. What does this say about a nonzero scalar CCR?
\end{exercise}
\noindent\emph{Solution.} The product rule gives $[A,B^n]=nB^{n-1}$. No power of $B$ vanishes, by this identity and $B^0=1$. Taking norms gives $n\|B^{n-1}\|\le2\|A\|\|B\|\|B^{n-1}\|$, impossible for all $n$. Rescaling one operator reduces any nonzero scalar commutator to this one. The conclusion concerns the fields, not their Weyl unitaries.

\begin{exercise}[Advanced: an isomorphism without a unitary]\label{f:ex-implementation}
In units $\hbar=1$, take the algebraic direct sum of countably many standard real label planes $(a_j,b_j)$ (each vector has finite support). Use the Gaussian state with $\mu(v,v)=\tfrac12\sum_j(a_j^2+b_j^2)$. Show that $S:(a_j,b_j)\mapsto(ja_j,b_j/j)$, $j\ge1$, induces a Weyl automorphism not unitarily implementable in this state's representation.
\end{exercise}
\noindent\emph{Solution.} The map and its inverse preserve finite support and each plane's alternating form. Take $v_n$ with only $a_n=1/n$ nonzero. Then $\|\pi(W(v_n))\Omega-\Omega\|^2=2(1-e^{-1/(4n^2)})\to0$. Each fixed $W(w)$ commutes with $W(v_n)$ for all sufficiently large $n$ by finite support. Convergence to the identity thus holds on a dense span of $\pi(W(w))\Omega$, and strongly everywhere by the unit norm bound. Unitary conjugation would give the same strong limit for $\pi(W(Sv_n))$, but its squared displacement of $\Omega$ is $2(1-e^{-1/4})$ for every $n$. The algebraic isomorphism has not failed; it is the proposed operator on this Hilbert space that does not exist.
